\documentclass[11pt]{article}
\usepackage[margin=1in]{geometry}
\usepackage[T1]{fontenc}
\usepackage[utf8]{inputenc}
\usepackage{hyperref}
\usepackage{longtable}
\usepackage{booktabs}
\setlength{\parindent}{0pt}
\setlength{\parskip}{0.7em}

\title{Phase-by-Phase Execution Plan\\Long Island Traffic Accident Cluster Discovery in R}
\author{AMS 597 Project Team}
\date{\today}

\begin{document}
\maketitle
\tableofcontents
\newpage

% -------------------------------------------------------
\section{Overview and Design Principles}

This document translates the SRS into a concrete, step-by-step execution plan for the Long Island accident clustering project. Each phase has a single, well-scoped objective, a precise list of tasks, explicit done criteria, and notes on known data issues that must be resolved within that phase.

The validated dataset (\texttt{li\_accidents\_clean.csv}) contains \textbf{33,375 rows} and \textbf{48 columns}. The analysis is split into two complementary pipelines:

\begin{enumerate}
  \item \textbf{Scenario pipeline}: mixed-type accident archetype discovery using \texttt{clustMixType::kproto()}
  \item \textbf{Spatial pipeline}: highway corridor hotspot discovery using \texttt{dbscan::dbscan()}
\end{enumerate}

Each pipeline proceeds independently and their results are joined for interpretation only in Phase~9.

\subsection{Confirmed Data Issues Requiring Action}

A pre-analysis audit of \texttt{li\_accidents\_clean.csv} revealed the following issues that must be resolved during execution:

\begin{longtable}{|p{4.5cm}|p{4.5cm}|p{5cm}|}
\hline
\textbf{Variable} & \textbf{Issue} & \textbf{Resolution} \\
\hline
\texttt{Duration\_Minutes} & Max = 2,236,406 min. 3,436 rows (10.3\%) above IQR upper fence of 270 min. 6 records exceed one week. & Cap at 99th percentile ($\approx$480 min) and apply log transform. Addressed in Phase~3. \\
\hline
\texttt{Distance(mi)} & 35.1\% of rows have Distance = 0. 251 records exceed 10 miles. Likely a reporting artifact, not a genuine trip distance. & Exclude from the clustering feature matrix. Addressed in Phase~4. \\
\hline
\texttt{Wind\_Direction} & String duplicates present: ``South'' and ``S'', ``West'' and ``W'', ``North'' and ``N'', ``East'' and ``E'' all coexist. & Normalize to 8-direction abbreviation scheme. Addressed in Phase~3. \\
\hline
\texttt{Bump}, \texttt{Roundabout}, \texttt{Turning\_Loop}, \texttt{Traffic\_Calming}, \texttt{No\_Exit} & All are 0\% TRUE across all 33,375 rows. Carry zero signal for clustering. & Drop from all feature matrices. Addressed in Phase~4. \\
\hline
\texttt{Give\_Way} & Only 0.9\% TRUE. Near-zero variance. & Exclude from core feature matrix. Addressed in Phase~4. \\
\hline
\texttt{Precipitation(in)} & 20.15\% missing. & Exclude from core model per SRS. Confirmed in Phase~2. \\
\hline
\texttt{Wind\_Chill(F)} & 15.01\% missing and highly correlated with Temperature. & Exclude from all models. Confirmed in Phase~2. \\
\hline
\texttt{End\_Lat / End\_Lng} & 39.21\% missing. & Exclude; use Start coordinates only. Confirmed in Phase~2. \\
\hline
\end{longtable}

% -------------------------------------------------------
\newpage
\section{Phase 1: Environment and Reproducibility Setup}

\subsection{Objective}

Establish a reproducible R environment before any data is loaded or transformed, so that all subsequent phases produce identical results across machines and sessions.

\subsection{Tasks}

\begin{enumerate}
  \item Install and load required packages. The full list is:
  \begin{itemize}
    \item \texttt{tidyverse} (data manipulation and plotting)
    \item \texttt{lubridate} (timestamp parsing)
    \item \texttt{clustMixType} (k-prototypes clustering)
    \item \texttt{dbscan} (density-based spatial clustering)
    \item \texttt{cluster} (Gower distance reference and silhouette)
    \item \texttt{fpc} (cluster stability helpers)
    \item \texttt{ggplot2} (visualization)
    \item \texttt{sf} or \texttt{geosphere} (coordinate projection for DBSCAN)
    \item \texttt{factoextra} (optional: validation plots)
  \end{itemize}

  \item Set a global random seed at the top of the master script:
  \begin{verbatim}
set.seed(42)
  \end{verbatim}
  All subsequent calls that accept a \texttt{set.seed} argument or \texttt{nstart} parameter must inherit or re-declare this seed to ensure determinism.

  \item Organize the working directory so that:
  \begin{itemize}
    \item \texttt{li\_accidents\_clean.csv} is the sole data input
    \item the R script is named \texttt{li\_clustering.R} (or \texttt{li\_clustering.Rmd})
    \item output tables and figures write to a \texttt{/output/} subfolder
    \item no intermediate \texttt{.RData} blobs are committed to the repository
  \end{itemize}

  \item Confirm that the R version and package versions are recorded. Add a short \texttt{sessionInfo()} call at the bottom of the script or Rmd for reproducibility.
\end{enumerate}

\subsection{Done Criteria}

\begin{itemize}
  \item All packages install without errors on the target machine.
  \item The master script runs from a clean R session without manual interventions.
  \item The folder structure matches the layout described above.
\end{itemize}

% -------------------------------------------------------
\newpage
\section{Phase 2: Data Ingestion and Raw Profiling}

\subsection{Objective}

Load the raw dataset, confirm its structure matches expectations from the SRS, and produce a documented profiling report before any transformation is applied. This phase is read-only with respect to the data.

\subsection{Tasks}

\begin{enumerate}
  \item Load the CSV:
  \begin{verbatim}
df_raw <- read.csv("li_accidents_clean.csv", stringsAsFactors = FALSE)
  \end{verbatim}

  \item Verify dimensions: expect 33,375 rows and 48 columns. Halt if either dimension is unexpected.

  \item Confirm the county split:
  \begin{itemize}
    \item Nassau: 19,390 rows
    \item Suffolk: 13,838 rows
    \item Queens: 147 rows (western LI edge, retained but not the focus)
  \end{itemize}

  \item Confirm the severity distribution:
  \begin{itemize}
    \item Severity 1: 161
    \item Severity 2: 28,357
    \item Severity 3: 3,938
    \item Severity 4: 919
  \end{itemize}
  Note: Severity 2 accounts for 85\% of records. Severity is not a clustering input.

  \item Compute and tabulate missingness for every column. Confirm the following critical values:
  \begin{itemize}
    \item \texttt{End\_Lat} / \texttt{End\_Lng}: $\approx$39.2\% missing
    \item \texttt{Precipitation(in)}: $\approx$20.1\% missing
    \item \texttt{Wind\_Chill(F)}: $\approx$15.0\% missing
    \item All other clustering candidates: below 1\% missing
  \end{itemize}

  \item Confirm that five near-zero binary flags carry no signal:
  \begin{itemize}
    \item \texttt{Bump}, \texttt{Roundabout}, \texttt{Turning\_Loop}, \texttt{Traffic\_Calming}, \texttt{No\_Exit}: all 0\% TRUE
  \end{itemize}

  \item Confirm the \texttt{Duration\_Minutes} anomaly:
  \begin{itemize}
    \item Median $\approx$ 78 min
    \item Mean $\approx$ 449 min (inflated by outliers)
    \item Max = 2,236,406 min
    \item 21 records exceed 1,440 min (one day)
    \item 6 records exceed 10,080 min (one week)
  \end{itemize}

  \item Confirm that \texttt{Distance(mi)} has 35.1\% zero values and a max of 35.6 miles.

  \item Inspect \texttt{Wind\_Direction} for string duplicates (``South'' vs ``S'', ``West'' vs ``W'', etc.).

  \item Tabulate \texttt{Weather\_Condition}: confirm that the top 6--8 labels account for the majority of rows and the long tail contains many sparse labels requiring grouping.

  \item Document all findings in a short profiling summary. This summary feeds directly into the data cleaning justifications in the final report.
\end{enumerate}

\subsection{Done Criteria}

\begin{itemize}
  \item All dimension, county, severity, and missingness checks pass.
  \item Known anomalies (Duration outliers, Distance zeros, Wind\_Direction duplicates) are confirmed and recorded.
  \item No data transformations have been applied yet.
\end{itemize}

% -------------------------------------------------------
\newpage
\section{Phase 3: Feature Engineering}

\subsection{Objective}

Derive analysis-ready features from raw columns. This phase produces all engineered variables that will be used in the clustering input matrices.

\subsection{Tasks}

\subsubsection*{3A. Timestamp Features}

\begin{enumerate}
  \item Parse \texttt{Start\_Time} to a \texttt{POSIXct} object using \texttt{lubridate::ymd\_hms()}.
  \item Derive the following variables:
  \begin{itemize}
    \item \texttt{hour}: integer 0--23
    \item \texttt{hour\_band}: factor with levels \texttt{Early\_AM} (0--5), \texttt{AM\_Rush} (6--9), \texttt{Midday} (10--14), \texttt{PM\_Rush} (15--18), \texttt{Evening} (19--21), \texttt{Late\_Night} (22--23)
    \item \texttt{weekday}: factor ordered Monday through Sunday
    \item \texttt{weekend}: logical TRUE if Saturday or Sunday
    \item \texttt{month}: integer 1--12
    \item \texttt{season}: factor with levels \texttt{Winter} (Dec--Feb), \texttt{Spring} (Mar--May), \texttt{Summer} (Jun--Aug), \texttt{Fall} (Sep--Nov)
    \item \texttt{rush\_hour}: logical TRUE if \texttt{hour} is in 6--9 or 15--18
  \end{itemize}
  \item Verify that no \texttt{NA} values appear in any derived temporal feature.
\end{enumerate}

\subsubsection*{3B. Duration Treatment}

\begin{enumerate}
  \item Confirm that the 99th percentile of \texttt{Duration\_Minutes} is approximately 480 minutes.
  \item Create \texttt{Duration\_cap}: replace any value above 480 with 480. This caps the 1\% of extreme outlier records without discarding them.
  \item Create \texttt{log\_Duration}: apply \texttt{log1p(Duration\_cap)} to compress the right-skewed distribution. This variable is a candidate for the scenario feature matrix.
  \item Do not delete the original \texttt{Duration\_Minutes} column; it is retained for reference.
  \item Verify: \texttt{summary(log\_Duration)} should show no extreme values and a reasonable distribution.
\end{enumerate}

\subsubsection*{3C. Weather Grouping}

\begin{enumerate}
  \item Count the number of distinct \texttt{Weather\_Condition} labels (approximately 50--60 in the raw data).
  \item Collapse labels into a \texttt{weather\_group} factor with the following levels:
  \begin{itemize}
    \item \texttt{Clear}: ``Fair'', ``Clear'', ``Fair / Windy''
    \item \texttt{Cloudy}: ``Mostly Cloudy'', ``Partly Cloudy'', ``Scattered Clouds'', ``Overcast'', ``Cloudy'', ``Mostly Cloudy / Windy'', ``Partly Cloudy / Windy'', ``Cloudy / Windy''
    \item \texttt{Rain}: ``Light Rain'', ``Rain'', ``Heavy Rain'', ``Light Drizzle'', ``Light Rain / Windy'', ``Heavy Rain / Windy'', ``Drizzle''
    \item \texttt{Snow\_Ice}: ``Light Snow'', ``Snow'', ``Heavy Snow'', ``Wintry Mix'', ``Light Freezing Rain'', ``Ice Pellets'', ``Sleet''
    \item \texttt{Fog\_Haze}: ``Fog'', ``Haze'', ``Smoke'', ``Mist'', ``Patches of Fog'', ``Shallow Fog''
    \item \texttt{Thunderstorm}: any label containing ``Thunder''
    \item \texttt{Other}: all remaining labels with fewer than 30 records
  \end{itemize}
  \item Verify the distribution of \texttt{weather\_group} with a frequency table. The \texttt{Clear} and \texttt{Cloudy} groups should dominate; \texttt{Other} should contain fewer than 200 records total.
\end{enumerate}

\subsubsection*{3D. Wind Direction Normalization}

\begin{enumerate}
  \item Identify all unique values of \texttt{Wind\_Direction}. The raw data contains both abbreviated and full-word forms: ``South'' and ``S'', ``West'' and ``W'', ``North'' and ``N'', ``East'' and ``E''.
  \item Normalize to standard 8-direction abbreviations: N, NE, E, SE, S, SW, W, NW.
  \item Map ``CALM'' and ``Variable'' / ``VAR'' to a separate level \texttt{Variable\_Calm}.
  \item Confirm the final \texttt{wind\_dir\_norm} factor has at most 9 levels.
  \item This variable is optional for the core scenario model; include it only if the number of factor levels does not dominate the Gower distance calculation. If in doubt, exclude it and note this decision.
\end{enumerate}

\subsubsection*{3E. Binary Flag Audit}

\begin{enumerate}
  \item Confirm zero-variance flags: \texttt{Bump}, \texttt{Roundabout}, \texttt{Turning\_Loop}, \texttt{Traffic\_Calming}, \texttt{No\_Exit} are all identically FALSE. Mark these as excluded.
  \item Confirm near-zero flags: \texttt{Give\_Way} (0.9\% TRUE), \texttt{Station} (0.7\% TRUE), \texttt{Railway} (0.15\% TRUE), \texttt{Amenity} (0.34\% TRUE). These carry minimal signal. Exclude \texttt{Give\_Way}, \texttt{Station}, \texttt{Railway}, and \texttt{Amenity} from the core scenario matrix. Document this decision.
  \item Retain for the core scenario matrix: \texttt{Crossing} (10.1\%), \texttt{Junction} (16.0\%), \texttt{Traffic\_Signal} (9.5\%), \texttt{Stop} (1.0\%).
  \item Convert all retained binary flags to \texttt{factor} type with levels \texttt{TRUE} and \texttt{FALSE} so \texttt{kproto()} treats them as categorical.
\end{enumerate}

\subsection{Done Criteria}

\begin{itemize}
  \item All engineered features are present in the working data frame with correct types.
  \item \texttt{weather\_group}, \texttt{hour\_band}, \texttt{season}, \texttt{weekday} are factors.
  \item \texttt{log\_Duration} exists and has no extreme values.
  \item \texttt{Wind\_Direction} is normalized to at most 9 levels.
  \item All zero-variance and near-zero-variance flags are marked for exclusion.
\end{itemize}

% -------------------------------------------------------
\newpage
\section{Phase 4: Build Clustering Input Matrices}

\subsection{Objective}

Construct the final feature matrices for each pipeline. Nothing outside these matrices enters the clustering algorithms.

\subsection{Tasks}

\subsubsection*{4A. Scenario Matrix}

\begin{enumerate}
  \item The scenario feature matrix (\texttt{df\_scenario}) contains exactly the following columns:

  \textit{Temporal factors:}
  \begin{itemize}
    \item \texttt{hour\_band} (factor)
    \item \texttt{weekday} (factor)
    \item \texttt{weekend} (factor or logical converted to factor)
    \item \texttt{season} (factor)
    \item \texttt{rush\_hour} (factor)
    \item \texttt{Sunrise\_Sunset} (factor: Day / Night)
  \end{itemize}

  \textit{Weather factors:}
  \begin{itemize}
    \item \texttt{weather\_group} (factor, 7 levels)
  \end{itemize}

  \textit{Numeric weather:}
  \begin{itemize}
    \item \texttt{Temperature(F)} (numeric)
    \item \texttt{Humidity(\%)} (numeric)
    \item \texttt{Pressure(in)} (numeric)
    \item \texttt{Visibility(mi)} (numeric)
    \item \texttt{Wind\_Speed(mph)} (numeric)
  \end{itemize}

  \textit{Road context factors:}
  \begin{itemize}
    \item \texttt{Crossing} (factor)
    \item \texttt{Junction} (factor)
    \item \texttt{Traffic\_Signal} (factor)
    \item \texttt{Stop} (factor)
  \end{itemize}

  \textit{Duration (optional):}
  \begin{itemize}
    \item \texttt{log\_Duration} (numeric) --- include only if less than 5\% of values are NA after capping
  \end{itemize}

  \item Confirm the following variables are \textbf{not} present in \texttt{df\_scenario}:
  \begin{itemize}
    \item \texttt{Severity}, \texttt{ID}, \texttt{Source}, \texttt{Description}, \texttt{Street}, \texttt{City}, \texttt{Zipcode}, \texttt{nzip}
    \item \texttt{Start\_Lat}, \texttt{Start\_Lng}, \texttt{End\_Lat}, \texttt{End\_Lng}
    \item \texttt{County} (used for post-cluster comparison only)
    \item \texttt{Distance(mi)} (35\% zeros, not a scenario characteristic)
    \item \texttt{Precipitation(in)}, \texttt{Wind\_Chill(F)} (excessive missingness)
    \item All zero/near-zero flags identified in Phase 3E
  \end{itemize}

  \item Handle remaining missing values in \texttt{df\_scenario}:
  \begin{itemize}
    \item For numeric columns with fewer than 1\% missing (Temperature, Humidity, Pressure, Visibility, Wind\_Speed): impute with the column median. Record how many rows were imputed.
    \item For factor columns: if fewer than 10 rows are missing, impute with the mode. If more, treat as a separate \texttt{Unknown} level.
  \end{itemize}

  \item Confirm the final \texttt{df\_scenario} has zero missing values before proceeding to Phase~5.

  \item Record the final column list and row count in a comment block at the top of the clustering script.
\end{enumerate}

\subsubsection*{4B. Spatial Matrix}

\begin{enumerate}
  \item The spatial feature matrix (\texttt{df\_spatial}) contains only:
  \begin{itemize}
    \item \texttt{x\_km}: Start\_Lng converted to approximate kilometers
    \item \texttt{y\_km}: Start\_Lat converted to approximate kilometers
  \end{itemize}

  \item Use the flat-earth approximation centered on the LI midpoint ($\approx 40.8^\circ$N, $-73.2^\circ$W):
  \begin{verbatim}
lat0 <- 40.8
x_km <- (df$Start_Lng - (-73.2)) * 111.32 * cos(lat0 * pi / 180)
y_km <- (df$Start_Lat - lat0) * 110.574
  \end{verbatim}
  This produces distances in kilometers for the DBSCAN \texttt{eps} parameter.

  \item Confirm there are no missing values in \texttt{Start\_Lat} or \texttt{Start\_Lng}.

  \item Confirm the spatial extent of the projected coordinates corresponds to Long Island's east-west span ($\approx$170 km) and north-south span ($\approx$40 km).
\end{enumerate}

\subsection{Done Criteria}

\begin{itemize}
  \item \texttt{df\_scenario} has zero NAs, correct column types, and the documented variable set.
  \item \texttt{df\_spatial} has two numeric columns, no NAs, and correct geographic extent.
  \item Both matrices have 33,375 rows (same row order as the original data for label joining).
\end{itemize}

% -------------------------------------------------------
\newpage
\section{Phase 5: k-Prototypes Scenario Clustering}

\subsection{Objective}

Run the k-prototypes algorithm on \texttt{df\_scenario} for candidate values of $k$, using Gower-style normalization, repeated starts, and a consistent random seed.

\subsection{Tasks}

\begin{enumerate}
  \item Verify that all factor columns in \texttt{df\_scenario} are \texttt{factor} type and all numeric columns are \texttt{numeric} type. \texttt{kproto()} will error if this is not satisfied.

  \item Set \texttt{lambda = NULL} to let \texttt{kproto()} estimate the numeric-to-categorical weight ratio automatically. Document the estimated lambda in the report.

  \item Use \texttt{type = "gower"} to range-normalize numeric variables before computing distances. This prevents high-variance numeric columns (e.g.\ Temperature) from dominating the distance calculation.

  \item Run \texttt{kproto()} for $k \in \{3, 4, 5, 6, 7, 8\}$ with \texttt{nstart = 10}. Store all six solutions in a named list:
  \begin{verbatim}
kp_results <- list()
for (k in 3:8) {
  set.seed(42)
  kp_results[[as.character(k)]] <- kproto(
    x     = df_scenario,
    k     = k,
    nstart = 10,
    type  = "gower",
    verbose = FALSE
  )
}
  \end{verbatim}

  \item For each solution, record:
  \begin{itemize}
    \item within-cluster total distance (the objective function value, \texttt{\$tot.withinss})
    \item cluster sizes (check for any cluster with fewer than 100 rows)
    \item whether any cluster is trivially dominated by a single factor level
  \end{itemize}

  \item Plot \texttt{tot.withinss} against $k$ to produce an elbow plot. Label each point with its $k$ value.

  \item Flag any $k$ solution where the smallest cluster contains fewer than 200 rows as potentially unstable. Do not automatically discard it, but note it for Phase~6 stability checks.
\end{enumerate}

\subsection{Done Criteria}

\begin{itemize}
  \item Six candidate \texttt{kproto} solutions exist in \texttt{kp\_results}.
  \item Elbow plot is produced and saved.
  \item Cluster sizes are tabulated for all six solutions.
  \item No solution crashes due to missing values or type mismatches.
\end{itemize}

% -------------------------------------------------------
\newpage
\section{Phase 6: Cluster Validation and Stability}

\subsection{Objective}

Select the best value of $k$ using internal validation metrics and bootstrap stability, then confirm that the chosen solution is academically defensible.

\subsection{Tasks}

\subsubsection*{6A. Internal Validation}

\begin{enumerate}
  \item Run \texttt{clustMixType::validation\_kproto()} across the candidate $k$ range. This function computes internal indices suited to mixed-type data (including within-cluster compactness relative to between-cluster separation):
  \begin{verbatim}
val_result <- validation_kproto(
  method = "kproto",
  object = df_scenario,
  k      = 3:8,
  nstart = 5,
  type   = "gower"
)
  \end{verbatim}

  \item Inspect and compare the validation indices across all $k$ values. Look for:
  \begin{itemize}
    \item a clear ``knee'' in the within-cluster distance
    \item consistently better separation scores at a particular $k$
    \item no sharp degradation in balance between $k$ and $k+1$
  \end{itemize}

  \item Do not choose $k$ solely to maximize one metric. Balance interpretability (can each cluster be described meaningfully?) against validation evidence.
\end{enumerate}

\subsubsection*{6B. Bootstrap Stability}

\begin{enumerate}
  \item Run \texttt{clustMixType::stability\_kproto()} on the top two candidate $k$ values from 6A:
  \begin{verbatim}
stab <- stability_kproto(
  object = df_scenario,
  k      = k_best,
  B      = 50,
  type   = "gower"
)
  \end{verbatim}

  \item Inspect the average Jaccard similarity for each cluster across bootstrap replicates. A Jaccard above 0.6 for all clusters is considered stable; above 0.75 is strong.

  \item If the chosen $k$ has one or more clusters with Jaccard below 0.5, note this as a limitation and prefer the next-best validated $k$.

  \item Record both the validation indices and the Jaccard summary in a table for the report.
\end{enumerate}

\subsubsection*{6C. Cluster Profile Check}

\begin{enumerate}
  \item For the selected $k$, run \texttt{clustMixType::clprofiles()} to produce factor-level bar plots and numeric boxplots for each cluster.

  \item Confirm that at least two clusters have a clearly different dominant weather group, hour band, or roadway context. If all clusters look nearly identical in their profiles, the chosen $k$ is too high or the feature matrix needs revision.

  \item Attach plain-English working names to the clusters at this stage. These are preliminary labels that will be refined in Phase~9 after severity and county profiling. Examples:
  \begin{itemize}
    \item \textit{AM Rush Corridor}: morning hour band, high Junction, mostly Clear weather
    \item \textit{Night Low-Visibility}: Night Sunrise\_Sunset, Fog or Rain weather group
    \item \textit{Weekend Surface Road}: weekend TRUE, Traffic\_Signal heavy, Midday or Evening hour band
    \item \textit{Winter Adverse}: Snow\_Ice weather group, low Temperature, low Visibility
  \end{itemize}
\end{enumerate}

\subsection{Done Criteria}

\begin{itemize}
  \item One final value of $k$ is selected and formally justified with validation and stability evidence.
  \item Each cluster has a working plain-English name.
  \item A cluster sizes table and the Jaccard summary are ready for the final report.
\end{itemize}

% -------------------------------------------------------
\newpage
\section{Phase 7: DBSCAN Spatial Hotspot Clustering}

\subsection{Objective}

Identify dense accident corridor clusters along Long Island's major east-west highways using density-based spatial clustering. This pipeline runs entirely on the projected coordinate matrix \texttt{df\_spatial}.

\subsection{Tasks}

\subsubsection*{7A. Parameter Selection}

\begin{enumerate}
  \item Run a k-nearest-neighbor distance plot to choose \texttt{eps}:
  \begin{verbatim}
knn_dist <- kNNdist(df_spatial, k = 5)
knn_sorted <- sort(knn_dist)
plot(knn_sorted, type = "l", xlab = "Points sorted by distance",
     ylab = "5-NN distance (km)", main = "kNN Distance Plot")
  \end{verbatim}

  \item Identify the ``elbow'' in the kNN distance curve. On Long Island's highway-dominated coordinate set, expect the elbow to fall in the range of 0.3--0.8 km given the density of highway segments.

  \item Set \texttt{minPts} between 10 and 30. A higher \texttt{minPts} suppresses small isolated clusters and emphasizes true highway-density zones. Start with \texttt{minPts = 15} as the initial candidate.

  \item Record the chosen \texttt{eps} and \texttt{minPts} values and justify them with reference to the kNN distance plot.
\end{enumerate}

\subsubsection*{7B. Run DBSCAN}

\begin{enumerate}
  \item Run DBSCAN with the selected parameters:
  \begin{verbatim}
db_result <- dbscan(df_spatial, eps = eps_chosen, minPts = minPts_chosen)
  \end{verbatim}

  \item Tabulate the results:
  \begin{itemize}
    \item number of non-noise clusters (cluster label $> 0$)
    \item size of each cluster
    \item percentage of rows labeled as noise (cluster label $= 0$)
  \end{itemize}

  \item Accept the solution if the noise percentage is between 5\% and 30\%. If noise exceeds 50\%, \texttt{eps} is too small. If there is only one cluster, \texttt{eps} is too large. Adjust and re-run if needed, documenting each trial.

  \item Inspect cluster shapes by plotting \texttt{x\_km} vs \texttt{y\_km} colored by cluster label. Expect elongated east-west shapes corresponding to the LIE ($\approx$latitude 40.77), Southern State Pkwy ($\approx$latitude 40.70), and Northern State Pkwy ($\approx$latitude 40.80).
\end{enumerate}

\subsubsection*{7C. Corridor Assignment}

\begin{enumerate}
  \item For each non-noise DBSCAN cluster, compute the centroid latitude and longitude. Cross-reference centroid latitude against known corridor latitudes to assign a corridor label:

  \begin{center}
  \begin{tabular}{|l|l|}
  \hline
  \textbf{Centroid Latitude (approx)} & \textbf{Corridor} \\
  \hline
  40.74 -- 40.76 & Long Island Expressway (I-495) \\
  40.68 -- 40.72 & Southern State Parkway \\
  40.78 -- 40.82 & Northern State Parkway \\
  40.64 -- 40.68 & Sunrise Highway \\
  40.72 -- 40.76 (north-south elongated) & Meadowbrook / Wantagh Pkwy \\
  \hline
  \end{tabular}
  \end{center}

  \item Clusters that do not align with any known major corridor are labeled as \textit{Surface Road Cluster} or by the dominant city name.

  \item Record which major highways appear as DBSCAN clusters and which do not. Note any unexpected hotspot zones.
\end{enumerate}

\subsection{Done Criteria}

\begin{itemize}
  \item DBSCAN solution has between 3 and 15 non-noise clusters with noise between 5\% and 30\%.
  \item Cluster shapes visually correspond to LI highway corridor alignments.
  \item Each cluster has a corridor label or a surface-road label.
  \item Parameter selection is justified with the kNN distance plot.
\end{itemize}

% -------------------------------------------------------
\newpage
\section{Phase 8: Post-Clustering Interpretation}

\subsection{Objective}

Join both sets of cluster labels back to the original data and produce the interpretation tables and profiles that will drive the final report.

\subsection{Tasks}

\subsubsection*{8A. Label Joining}

\begin{enumerate}
  \item Add scenario cluster labels to the original data frame:
  \begin{verbatim}
df_raw$scenario_cluster <- kp_final$cluster
  \end{verbatim}

  \item Add spatial cluster labels:
  \begin{verbatim}
df_raw$spatial_cluster <- db_result$cluster
  \end{verbatim}

  \item Confirm that both label vectors have 33,375 elements and no NAs.
\end{enumerate}

\subsubsection*{8B. Severity Profiling by Scenario Cluster}

\begin{enumerate}
  \item Compute the percentage of Severity 3 and Severity 4 accidents within each scenario cluster.
  \item Produce a summary table:

  \begin{center}
  \begin{tabular}{|l|r|r|r|r|r|}
  \hline
  \textbf{Cluster} & \textbf{N} & \textbf{Sev1 \%} & \textbf{Sev2 \%} & \textbf{Sev3 \%} & \textbf{Sev4 \%} \\
  \hline
  (fill at runtime) & & & & & \\
  \hline
  \end{tabular}
  \end{center}

  \item Flag any cluster with a Severity 4 share more than twice the dataset average (2.75\%) as a high-risk archetype.
\end{enumerate}

\subsubsection*{8C. County Composition by Scenario Cluster}

\begin{enumerate}
  \item Compute the Nassau and Suffolk share within each scenario cluster.
  \item Produce a cross-tabulation. Identify whether any cluster is predominantly Nassau or predominantly Suffolk, and interpret this in light of the highway geography.
\end{enumerate}

\subsubsection*{8D. Weather and Time Profile by Cluster}

\begin{enumerate}
  \item For each scenario cluster, compute:
  \begin{itemize}
    \item dominant \texttt{weather\_group} (top label and its share)
    \item dominant \texttt{hour\_band}
    \item Day vs Night share (\texttt{Sunrise\_Sunset})
    \item weekend share
    \item mean Temperature, Visibility, and Wind\_Speed
  \end{itemize}
  \item Produce a compact cluster profile table for the report.
\end{enumerate}

\subsubsection*{8E. Spatial--Scenario Cross-Tabulation}

\begin{enumerate}
  \item For non-noise DBSCAN records only, cross-tabulate \texttt{spatial\_cluster} against \texttt{scenario\_cluster}.
  \item Identify whether corridor-specific hotspot zones are dominated by particular scenario archetypes. For example, does the LIE corridor cluster map predominantly to an AM Rush archetype? Does the surface road cluster map to a junction-heavy scenario?
  \item This cross-tabulation is a key insight of the two-pipeline design.
\end{enumerate}

\subsubsection*{8F. Final Cluster Naming}

\begin{enumerate}
  \item Using the full profile from 8B through 8E, finalize the plain-English name for each scenario cluster and each spatial hotspot cluster.
  \item Scenario cluster names should reflect the dominant time, weather, and road context. Examples:
  \begin{itemize}
    \item \textit{AM Rush --- Clear Highway}
    \item \textit{PM Merge Zone --- Congested Commute}
    \item \textit{Night Low-Visibility --- Fog and Rain}
    \item \textit{Winter Adverse --- Snow and Ice}
    \item \textit{Weekend Surface Road --- Signal and Crossing}
  \end{itemize}
  \item Spatial hotspot names should reference the corridor. Example: \textit{LIE Westbound Corridor}, \textit{Southern State Pkwy Eastbound Corridor}.
\end{enumerate}

\subsection{Done Criteria}

\begin{itemize}
  \item Severity, county, weather, and temporal profile tables are complete.
  \item Spatial--scenario cross-tabulation is produced.
  \item All clusters have finalized plain-English names.
  \item No analytical question from the SRS research questions is left unaddressed.
\end{itemize}

% -------------------------------------------------------
\newpage
\section{Phase 9: Export and Visualization}

\subsection{Objective}

Produce all report-ready tables, figures, and the final labeled dataset for submission.

\subsection{Tasks}

\begin{enumerate}
  \item \textbf{Elbow plot}: \texttt{tot.withinss} vs $k$ for the k-prototypes solutions (Phase~5).

  \item \textbf{Validation index plot}: internal validation metrics across $k = 3:8$ (Phase~6A).

  \item \textbf{Stability plot}: Jaccard bootstrap summary for the selected $k$ (Phase~6B).

  \item \textbf{Cluster profile plot}: \texttt{clprofiles()} output for the selected scenario clustering (Phase~6C). One panel per cluster showing factor bar charts and numeric boxplots.

  \item \textbf{kNN distance plot}: the 5-NN distance curve used to select DBSCAN \texttt{eps} (Phase~7A).

  \item \textbf{Spatial hotspot map}: scatter plot of \texttt{x\_km} vs \texttt{y\_km} colored by DBSCAN cluster label (Phase~7B). Noise points shown in grey. Highway corridor labels overlaid.

  \item \textbf{Severity-by-cluster table}: formatted table for the report showing each scenario cluster's Severity 3+4 share (Phase~8B).

  \item \textbf{County-by-cluster table}: Nassau vs Suffolk composition by scenario cluster (Phase~8C).

  \item \textbf{Spatial--scenario cross-tabulation}: heat-table or grouped bar chart (Phase~8E).

  \item \textbf{Labeled CSV export}: write the final data frame with \texttt{scenario\_cluster} and \texttt{spatial\_cluster} columns to \texttt{output/li\_accidents\_clustered.csv}.

  \item All figures should be saved at 150 DPI or higher in PNG or PDF format to \texttt{output/figures/}.

  \item All tables should be saved as CSV to \texttt{output/tables/} for easy insertion into the LaTeX report.
\end{enumerate}

\subsection{Done Criteria}

\begin{itemize}
  \item All nine figures exist in \texttt{output/figures/}.
  \item All tables exist in \texttt{output/tables/}.
  \item The labeled CSV is exported.
  \item Every visual is readable at print resolution without axis label truncation.
\end{itemize}

% -------------------------------------------------------
\newpage
\section{Phase 10: Final Report and Submission}

\subsection{Objective}

Assemble the final report, compile the PDF, and prepare the branch for review.

\subsection{Tasks}

\begin{enumerate}
  \item Write the final report in LaTeX. Sections to include:
  \begin{itemize}
    \item Introduction and motivation (derived from SRS Sections 1--3)
    \item Dataset description and data audit findings (Phase~2 summary)
    \item Feature engineering decisions (Phase~3, with justification table)
    \item Method selection rationale (derived from SRS Section~4)
    \item Scenario clustering results: elbow plot, chosen $k$, validation and stability, cluster profiles
    \item Spatial hotspot results: kNN plot, DBSCAN map, corridor assignments
    \item Interpretation: severity table, county table, spatial--scenario cross-tabulation
    \item Limitations: Duration cap, Precipitation exclusion, Distance exclusion, Severity 2 dominance
    \item Conclusion: which clusters are most relevant for safety attention
  \end{itemize}

  \item Keep the limitations section honest. In particular:
  \begin{itemize}
    \item acknowledge that Duration capping at 480 min discards real variability in the top 1\% of records
    \item acknowledge that Precipitation exclusion (20\% missing) may underrepresent rain-related accident patterns
    \item acknowledge that spatial hotspot shapes are sensitive to the choice of \texttt{eps}
  \end{itemize}

  \item Compile the PDF from the LaTeX source and verify all figures render correctly.

  \item Commit the following files to the \texttt{topic-3-nyc-accident-clustering} branch:
  \begin{itemize}
    \item \texttt{NYC\_Clustering\_SRS.tex} and \texttt{NYC\_Clustering\_SRS.pdf}
    \item \texttt{NYC\_Clustering\_Execution\_Plan.tex} and PDF
    \item \texttt{li\_clustering.R} (or \texttt{li\_clustering.Rmd})
    \item \texttt{output/} subfolder with all figures, tables, and the labeled CSV
    \item \texttt{li\_accidents\_clean.csv} (if not already committed)
  \end{itemize}

  \item Push the branch to the remote repository once the repository owner has added the contributor.
\end{enumerate}

\subsection{Done Criteria}

\begin{itemize}
  \item The final report PDF compiles without errors.
  \item All figures are embedded and readable.
  \item The branch is clean and the commit message describes the submission state clearly.
  \item The work addresses all five SRS research questions.
\end{itemize}

% -------------------------------------------------------
\newpage
\section{Phase Summary Table}

\begin{longtable}{|p{0.5cm}|p{3cm}|p{6.5cm}|p{4.5cm}|}
\hline
\textbf{\#} & \textbf{Phase} & \textbf{Key Output} & \textbf{Done Signal} \\
\hline
1 & Environment Setup & Clean R project, packages installed, seed set & Script runs from cold start \\
\hline
2 & Data Profiling & Missingness table, anomaly confirmation & All expected row/column counts verified \\
\hline
3 & Feature Engineering & \texttt{hour\_band}, \texttt{season}, \texttt{weather\_group}, \texttt{log\_Duration}, wind normalization & Zero NAs in all engineered features \\
\hline
4 & Build Input Matrices & \texttt{df\_scenario} (no leakage), \texttt{df\_spatial} (km coords) & Both matrices have 33,375 rows, zero NAs \\
\hline
5 & k-Prototypes Clustering & Six candidate solutions, elbow plot & No crashes; cluster sizes tabulated \\
\hline
6 & Validation \& Stability & Selected $k$, Jaccard summary, cluster profiles & One defensible solution with working names \\
\hline
7 & DBSCAN Spatial & Corridor hotspot clusters, kNN plot, map & 3--15 clusters, 5--30\% noise, corridor shapes confirmed \\
\hline
8 & Interpretation & Severity/county/weather tables, spatial--scenario cross-tab, final names & All SRS research questions answered \\
\hline
9 & Export & Figures, tables, labeled CSV & All outputs in \texttt{output/} folder \\
\hline
10 & Report \& Submission & Final LaTeX report, PDF, committed branch & PDF compiles, branch is clean \\
\hline
\end{longtable}

% -------------------------------------------------------
\section{Variable Decision Reference}

The following table provides a single-source-of-truth for every variable decision made in this execution plan.

\begin{longtable}{|p{4cm}|p{2.5cm}|p{7cm}|}
\hline
\textbf{Variable} & \textbf{Decision} & \textbf{Reason} \\
\hline
\texttt{hour\_band} & Include & Key temporal feature, derived from Start\_Time \\
\hline
\texttt{weekday} & Include & Commuter vs non-commuter pattern signal \\
\hline
\texttt{weekend} & Include & Captures Saturday/Sunday off-peak behavior \\
\hline
\texttt{season} & Include & Weather-linked seasonal variation \\
\hline
\texttt{rush\_hour} & Include & Direct commuter proxy \\
\hline
\texttt{Sunrise\_Sunset} & Include & Day/Night is a strong accident context variable \\
\hline
\texttt{weather\_group} & Include & Collapsed Weather\_Condition; reduces sparsity \\
\hline
\texttt{Temperature(F)} & Include & Numeric, no missing, strong season proxy \\
\hline
\texttt{Humidity(\%)} & Include & Numeric, no missing \\
\hline
\texttt{Pressure(in)} & Include & Numeric; 1.8\% mild IQR outliers are physically valid \\
\hline
\texttt{Visibility(mi)} & Include & Key safety variable; low-visibility records are real events \\
\hline
\texttt{Wind\_Speed(mph)} & Include & Numeric; 2.4\% high outliers are physically plausible \\
\hline
\texttt{Crossing} & Include & 10.1\% TRUE; meaningful road context \\
\hline
\texttt{Junction} & Include & 16.0\% TRUE; strongest road context flag \\
\hline
\texttt{Traffic\_Signal} & Include & 9.5\% TRUE; distinguishes surface roads from highways \\
\hline
\texttt{Stop} & Include & 1.0\% TRUE; low but non-trivial; kept as marginal signal \\
\hline
\texttt{log\_Duration} & Include (conditional) & After cap at 480 min; carries incident severity signal \\
\hline
\texttt{Precipitation(in)} & Exclude & 20.1\% missing; too high for routine imputation \\
\hline
\texttt{Wind\_Chill(F)} & Exclude & 15.0\% missing; redundant with Temperature \\
\hline
\texttt{Distance(mi)} & Exclude & 35.1\% zeros; reporting artifact, not scenario signal \\
\hline
\texttt{Amenity} & Exclude & 0.34\% TRUE; near-zero variance \\
\hline
\texttt{Give\_Way} & Exclude & 0.9\% TRUE; near-zero variance \\
\hline
\texttt{Station} & Exclude & 0.7\% TRUE; near-zero variance \\
\hline
\texttt{Railway} & Exclude & 0.15\% TRUE; near-zero variance \\
\hline
\texttt{Bump} & Exclude & 0.0\% TRUE; zero variance \\
\hline
\texttt{Roundabout} & Exclude & 0.0\% TRUE; zero variance \\
\hline
\texttt{Turning\_Loop} & Exclude & 0.0\% TRUE; zero variance \\
\hline
\texttt{Traffic\_Calming} & Exclude & 0.0\% TRUE; zero variance \\
\hline
\texttt{No\_Exit} & Exclude & 0.0\% TRUE; zero variance \\
\hline
\texttt{Severity} & Post-cluster only & Leakage risk; profile after clustering \\
\hline
\texttt{County} & Post-cluster only & Use for Nassau/Suffolk comparison after clustering \\
\hline
\texttt{Start\_Lat / Start\_Lng} & Spatial only & Used only in df\_spatial for DBSCAN \\
\hline
\texttt{End\_Lat / End\_Lng} & Exclude & 39.2\% missing \\
\hline
\texttt{ID, Source, Description, Street, City, Zipcode, nzip} & Exclude & Identifiers and free text; no clustering value \\
\hline
\end{longtable}

\end{document}
